\documentclass[a4paper]{article}

\usepackage{graphicx}
\usepackage{amsmath,amssymb}
\usepackage{minted}
\usepackage{multirow}
\usepackage{float}
\usepackage{todonotes}
\usepackage{forest}
\usepackage{xstring}
\usepackage{xcolor}
\usepackage[most]{tcolorbox}
\usepackage{xparse}
\usepackage{natbib}

\usetikzlibrary{graphs,graphdrawing}
\usegdlibrary{layered}

% for external graphviz
\input{/home/donkarlo/Dropbox/repo/nd_language_ai_project/src/nd_language_ai/formal/computer/programming/tex/latex/code_clips/external_graphviz.tex}

% for tags
\input{/home/donkarlo/Dropbox/repo/nd_language_ai_project/src/nd_language_ai/formal/computer/programming/tex/latex/parts/control_sequence/kind/semtag/semtag.tex}

% for links
\input{/home/donkarlo/Dropbox/repo/nd_language_ai_project/src/nd_language_ai/formal/computer/programming/tex/latex/code_clips/seehere.tex}

\title{}
\date{}
\author{Mohammad Rahmani}

\begin{document}
    \maketitle
    Ich kann nicht am Samstag für dich arbeiten, erstens weil ich kein Sklave bin und zweitens weil am Samstag ich sehr beschäftigt bin.
    
    Ich denke wirklich, dass du höflich bitten lernen musst, sonst macht niemand etwas für dich.
    
    Wenn du mir so geschrieben hast, dann, selbst wenn ich dir helfen wollen würde, würde ich es nicht machen, um dir klarzumachen, dass ich nicht dein Skalve bin.
    
    Ich weiß, dass das(Neutral ist benutzt worden, weil weil wir ein pronom brauchen der wir das Genus nicht sicher sind ), was du willst, sehr wichtig ist, aber falls ich dir dieses(demonstrativ Pronom, -s ist für Akkusativ) Mal(neutral) helfen würde, würdest du jedes Mal(n) denken, dass du mit so einem Ton(m), die Freunde immer wie Mitarbeiter insturieren   kanst.
    
    Endlich muss ich es nennen, obwohl ich nicht dein Sklave bin, kann ich dich auch nicht allein lassen, um vermeide dich in Leiden zu sehen. Deshalb werde ich anwesend sein, um dich zu helfen.



    
    
    
    
    
    \bibliographystyle{plainnat}
    \bibliography{/home/donkarlo/Dropbox/repo/nd_shared_research/bibliography/bibliography.bib}
\end{document}